\chapter{あとがき}

本文と関係ないことで無駄に1ページを費すことに対する反省の弁。


\clearpage
\thispagestyle{empty}
% ここから下は、かなり雑にレイアウトしてるので適宜修正すべし。

\begin{flushleft}
{\LARGE\faIcon{pen-to-square}}\quad{\large\sffamily 著者紹介}

\vspace{0.5\baselineskip}
{\sffamily{}何野誰某}
\vspace{0.3\baselineskip}

そのへんのおっさん。
\end{flushleft}

\vspace{1\baselineskip}% QRコードの周囲にほとんど空白が空かないので手で調整する
\begin{tblr}{cc}
\qrcode[height=2cm]{https://x.com/example_account} &
\qrcode[height=2cm]{https://example.com} \\
\href{https://x.com/example_account}{\faIcon{square-x-twitter} \ttfamily@example\_account} &
\url{https://example.com}
\end{tblr}
\vspace{1\baselineskip}



% 奥付は偶数ページに載せる
% 以下のコメント部分を生かすと自動で判定されるが、手作業で調整してもいいんじゃないかな
% (偶数ページなら改ページせず著者紹介の下のスペースに、奇数ページなら改ページ)
%\makeatletter
%\ifodd\c@CurrentPage
%  \clearpage\thispagestyle{empty}
%\fi
%\makeatother

\null\vfill

% 奥付本体
\begin{center}
\parbox{25\zw}{
{\large\sffamily{}かっこいいタイトル}
\vspace*{0.5\zh}
\hrule height 1pt % 太さ1pt横幅いっぱいの線
\vspace*{1\zh}
YYYY年MM月DD日 初版第1刷
\vspace*{1\zh}
\begin{center}
\begin{tblr}{ll}
著者 & 何野誰某 \\
発行 & Foobar Publisher \\
e-mail & \texttt{mail@address} \\
印刷 & ほげほげ印刷 \\
\end{tblr}
\end{center}
\vspace*{1\zh}
\hrule height 0.4pt
\null\hfill{\ttfamily rev. \revid}
}% end \parbox
\end{center}
